% Rédaction : KOUMTOUBRE Davy ; relecture : BADO Lionel
\chapter{Limites et perspectives}\label{chap:limites}

Le projet démontre une chaîne intégrée de traitement et d'exposition de transactions fictives.
Il s'agit d'un livrable pédagogique, pas d'un système prêt à traiter des opérations bancaires
réelles. Les limites suivantes doivent être prises en compte.

\section{Périmètre fonctionnel et données}
Les données sont entièrement synthétiques et produites par un générateur déterministe. Elles
permettent de reproduire les tests, mais ne représentent ni la diversité ni les contraintes de
confidentialité de données bancaires réelles. L'API ne réalise pas de mouvement de fonds : elle
enregistre et expose des objets de transaction.

Le workflow KNIME est conçu pour une exécution sur poste de travail et une production de fichiers
CSV. Cette étape reste distincte de l'API et de Docker. Pour une exploitation régulière, il faudrait
automatiser son exécution, superviser les fichiers produits et définir une stratégie de reprise en
cas d'échec de l'ETL.

\section{Sécurité et exploitation}
L'API ne comporte pas d'authentification, d'autorisation par rôle ni de limitation de débit. Les
services et identifiants fournis par défaut sont destinés au développement local. Avant toute
exposition sur un réseau non maîtrisé, il faudrait ajouter une gestion d'identité (par exemple une
clé d'API ou des jetons), protéger les secrets, configurer TLS et définir une politique d'accès
minimale à PostgreSQL.

Prometheus collecte les métriques et Grafana les présente, mais les règles d'alerte ne disposent
pas d'un service de notification tel qu'Alertmanager. Les journaux sont consultables dans les
sorties des conteneurs; ils ne sont pas centralisés ni conservés dans un système de recherche.
La pile Docker Compose vise un seul poste et ne fournit ni haute disponibilité, ni sauvegarde
automatisée, ni stratégie de déploiement multi-instance.

\section{Évolutions possibles}
Les suites naturelles sont l'automatisation du workflow KNIME, l'ajout d'une authentification et
d'une autorisation adaptées, la gestion centralisée des secrets, les sauvegardes/restaurations
testées, ainsi que l'envoi des alertes et la centralisation des journaux. Une mise en production
nécessiterait également des tests de charge plus représentatifs, une politique de rétention des
métriques et journaux, et une revue de sécurité de l'ensemble de la chaîne.